\subsection{Definition of a convex set}

For any two points x, y of an affine space E, the set of points \(\lambda x + \mu y\) where \(\lambda \geqslant 0\) , \(\mu \geqslant 0\) , \(\lambda + \mu = 1\) is called the closed segment with end points x and y; it reduces to a point when x = y. The complement of x in this segment is called the segment with end points x, y which is open at x and closed at y; it is empty if x = y. Finally the complement of \(\{x, y\}\) in the closed segment with end points x, y is called the open segment with end points x, y; it is empty when x = y.

DEFINITION 1. --- A subset A of an affine space E is convex if, for every two points x, y of A, the closed segment with end points x, y is contained in A.

As \((1 - \lambda)a + \lambda x = a + \lambda(x - a)\) , this definition is equivalent to the following: the set A is convex if, for every point \(a \in A\) , the transform of A by a homothety of centre a and ratio \(\lambda\) where \(0 < \lambda < 1\) , is contained in A (in other words, A is stable for these homotheties).

Examples. --- 1) Every linear affine variety of E (and in particular the empty set) is convex.

\begin{enumerate}
\def\labelenumi{\arabic{enumi})}
\setcounter{enumi}{1}
\item
  The only non-empty convex sets in \(\mathbf{R}\) are the intervals (GT, IV, § 2.4, prop. 1).
\item
  Let E be a vector space and \(\|x\|\) a norm on E; the unit ball B, formed by the points x such that \(\| x\| \leqslant 1\) , is convex since the relations \(\| x\| \leqslant 1\) , \(\| y\| \leqslant 1\) , imply for \(0\leqslant \lambda \leqslant 1\) that
\end{enumerate}

\[
\left\| \lambda x + (1 - \lambda) y \right\| \leqslant \lambda \| x \| + (1 - \lambda) \| y \| \leqslant \lambda + (1 - \lambda) = 1.
\]

Remark. --- Let A be a convex subset of a vector space E; for any scalars \(\alpha > 0\) and \(\beta > 0\) we have \(\alpha A + \beta A = (\alpha + \beta) A\) . In other words, for any \(x \in A\) , \(y \in A\) , there exists \(z \in A\) such that \((\alpha + \beta) z = \alpha x + \beta y\) ; in fact this relation can be written \(z = \frac{\alpha}{\alpha + \beta} x + \frac{\beta}{\alpha + \beta} y\) and we have \(\frac{\alpha}{\alpha + \beta} > 0\) , \(\frac{\beta}{\alpha + \beta} > 0\) and \(\frac{\alpha}{\alpha + \beta} + \frac{\beta}{\alpha + \beta} = 1\) , from which the assertion follows, on using def. 1.

PROPOSITION 1. --- Let \((x_{\iota})\) be a family of points of a convex subset A; every barycentre \(\sum_{\iota} \lambda_{\iota} x_{\iota}\) of the \(x_{\iota}\) formed using positive masses \(\lambda_{\iota}\) (such that \(\sum_{\iota} \lambda_{\iota} = 1\) and \(\lambda_{\iota} = 0\) except for finitely many of the indices, cf.~A, II, §9.3) belongs to A.

Clearly we need only consider the case when the indices are 1, 2, \ldots, p and \(\lambda_{i} > 0\) for each i; the proposition is trivial if p = 1; we prove the result by induction on p.~Put \(\mu = \sum_{i=1}^{p-1} \lambda_{i} > 0\) , and \(y = \sum_{i=1}^{p-1} \frac{\lambda_{i}}{\mu} x_{i}\) ; the induction hypothesis implies that \(y \in A\) . Now as \(\lambda_{p} = 1 - \mu\) and \(\sum_{i=1}^{p} \lambda_{i} x_{i} = \mu y + (1 - \mu) x_{p}\) , its follows from def. 1 that \(\sum_{i=1}^{p} \lambda_{i} x_{i}\) belongs to A.

PROPOSITION 2. --- Let E and F be two affine spaces and f be an affine linear mapping of E in F; then the image of a convex subset of E under f, and the inverse image of a convex subset of F under f are both convex.

The image under f of the closed segment with end points x, y is the closed segment with end points \(f(x)\) , \(f(y)\) , hence the first statement. We deduce that the inverse image of a closed segment of F under f contains each closed segment whose end points belong to it; the second statement of prop. 2 follows.

In particular the image of a convex set under a homothety or a translation is a convex set.

PROPOSITION 3. --- In the affine space E, let H be a hyperplane defined by the relation \(g(x) = 0\) , where g is a non-constant affine function on E. Then the half-spaces defined by the relations \(g(x) \geqslant 0\) , \(g(x) \leqslant 0\) , \(g(x) > 0\) , \(g(x) < 0\) are convex.

For these are the inverse images under g of intervals of R and thus are convex.

With the notations of prop. 3 the points of a subset M of an affine space are on the same side (resp. strictly on the same side) of the hyperplane H if M is contained in one of the half-spaces defined by \(g(x) \geqslant 0\) , \(g(x) \leqslant 0\) (resp. \(g(x) > 0\) or \(g(x) < 0\) ).

PROPOSITION 4. --- The points of A, a convex subset of an affine space E are strictly on the same side of a hyperplane H if, and only if, A does not meet H.

Clearly the condition is necessary. Conversely suppose that it is satisfied and let \(g(x) = 0\) , be an equation defining H (g is an affine linear mapping of E in R). The set \(g(\mathbf{A})\) is convex in R, therefore it is an interval, and \(0 \notin g(\mathbf{A})\) . Hence \(g(x)\) is of fixed sign for all \(x \in A\) .

